\begin{savequote}[8cm]
Wenn du es nicht einfach erklären kannst, hast du es nicht gut genug verstanden."

  \qauthor{--- Albert Einstein}
\end{savequote}

\newpage
% Analysis of results and conclusions, discussion of plans for future work
%\vspace{-4mm}
\section{Conclusion and Outlook} \label{conclusion} %\vspace{-3mm}
This report illustrates the first experimental study of quantum decoherence using accelerator-born neutrinos. As mentioned and referenced in earlier sections, phenomenological quantum decoherence studies as well as studies on quantum decoherence using astrophysical or atmospheric neutrinos can be found in the literature. However, this is the first experimental study of quantum decoherence considering accelerator-neutrinos. The special feature of this study, as opposed to a phenomenological study, is that this experimental analysis includes the flux, detector and neutrino interaction simulations with its statistical and systematic uncertainties as they are considered in the \acrshort{dune} and \acrshort{t2k} experiments. \\

In this report, I have shown the effects of decoherence with respect to particular parameter choices using the official neutrino oscillation analyses workflow in the \acrshort{mach3} framework in \acrshort{dune} and \acrshort{t2k}. After understanding the content and formalism of the vast topic of quantum decoherence, the essence of my work was to perform a first feasibility study (\cref{fig:deimos_plots}) and thereafter implement the \acrshort{nusquids} code into NuOscillator in order to freely use it as part of \acrshort{mach3}. I did this in both the \acrshort{mach3} code versions used in \acrshort{dune} and \acrshort{t2k}. After validating these implementations, I generated all plots shown (figs. \ref{fig:event_spectra_dune} to \ref{fig:nu_antinu_comp}) using the newly available oscillation engine and advertised this contribution to the official neutrino oscillation analyses in both experimental collaborations. As the \acrshort{nusquids} code is more powerful than just enabling studies on quantum decoherence studies, but also on other \acrshort{bsm} models such as \acrshort{liv}, \acrshort{nsi} and sterile neutrinos, this implementation enables collaborators to more quickly initiate new \acrshort{bsm} studies, allowing them to reach the core of a new analysis more efficiently and with reduced startup time. \\

Future investigations will choose a different reference energy $E_0$, for example $E_0 = 1$ MeV as the choice of $1$ GeV is around the most frequent neutrino energies in \acrshort{dune} and \acrshort{t2k} that are considered in this study. This will solely shift the decoherence strength $\Gamma_0$ by a factor of $\left(\frac{E_{0(old)}}{E_{0(new)}}\right)^n$ and not the physics under investigation. As it can be seen in most plots shown in this report, especially in \cref{fig:qd_t2k_gamma_varied_highlights1}, the artificial pivot point can be clearly seen at $E_0=1$ GeV and hence produces a less convenient point to consider when interpreting the results. \\

All plots shown show the reconstructed neutrino energy and not the true neutrino energy. Event rates in the true neutrino energy would enhance the signal seen from quantum decoherence.

The production of neutral current samples in \acrshort{t2k} is currently underway and will be included in future analyses. \\

In order to answer the central question of whether quantum decoherence can be observed in the \acrshort{dune} and \acrshort{t2k} experiments, a sensitivity study will be performed to obtain an estimate on the sensitivity of measuring the decoherence strength for different choices of the energy dependence with respect to these two experiments. Subsequently, I will estimate confidence level upper limits with respect to different choices of the energy scale. \acrshort{mach3} is a powerful tool to perform such statistical studies. One immediate next step is to consider likelihood scans of the parameters in question. As it was outlined in this report, there is a plethora of possibilities inherited in the open quantum systems formalism, in particular in the formulation of the decoherence matrix giving rise to a particular phenomenological model. Further models and alternative physics interpretations will be investigated in future analyses.